\documentclass[12pt, letterpaper]{article}

% --- Packages ---
\usepackage{geometry}
 \geometry{margin=1in}
\usepackage{graphicx}
\usepackage{amsmath, amssymb}
\usepackage{booktabs}
\usepackage{hyperref}
\usepackage{cite}
\usepackage{setspace}
\onehalfspacing
\usepackage{fontspec}
\usepackage{polyglossia}
\usepackage{float}
\usepackage{svg}
\usepackage{enumitem}
\setmainlanguage{vietnamese}

\usepackage{listings}
\usepackage{xcolor}

\lstset{
    language=C,
    basicstyle=\ttfamily\small,
    keywordstyle=\color{blue},
    stringstyle=\color{red},
    commentstyle=\color{green!60!black},
    breaklines=true,
    frame=single, % Tạo khung xung quanh code
    backgroundcolor=\color{gray!5}
}

\renewcommand{\lstlistingname}{Mã nguồn}

% --- Metadata ---
\title{Giải bài tập slide trang 29, 34, 35 Bài 5}
\author{Đỗ Hoàng Gia Bảo - 23020783 \\ \small Khoa Điện tử - Viễn Thông}
\date{\today}

\begin{document}

\maketitle

\begin{abstract}
Báo cáo sẽ trình bày các bài tập tại slide trang 29, 34, 35 \texttt{Bài-5}
\end{abstract}
% --- Front Matter (Roman Numerals) ---
\newpage
\tableofcontents

% --- Main Content (Arabic Numerals) ---
\newpage
\pagenumbering{arabic}
\setcounter{page}{3} % This automatically resets the counter to 1

\section{Bài tập slide trang 29}
\subsection{Nội dung bài tập}
\begin{table}[H]
    \centering
    \begin{tabular}{|c|c|c|c|}
        \toprule
        & \textbf{Allocation} & \textbf{Need} & \textbf{Available} \\
        & A B C & A B C & A B C \\
        \midrule
        $P_0$ & 0 1 0 & 7 4 3 & 2 3 0 \\
        $P_1$ & 3 0 2 & 0 2 0 &\\
        $P_2$ & 3 0 2 & 6 0 0 &\\
        $P_3$ & 2 1 1 & 0 1 1 &\\
        $P_4$ & 0 0 2 & 4 3 1 &\\
        \bottomrule
    \end{tabular}
\end{table}

\begin{itemize}
    \item Thực thi thuật toán an toàn thì chuỗi < $P_1, P_3, P_4, P_0, P_2$ > thoả mãn điều kiện an toàn
    \item Liệu yêu cầu (3, 3, 0) bởi $P_4$ được chấp nhận?
    \item Liệu yêu cầu (0, 2, 0) bởi $P_0$ được chấp nhận?
\end{itemize}

\subsection{Giải bài tập}
\subsubsection{Thực thi thuật toán an toàn}
\textbf{Khởi tạo 2 mảng Finish và Work:}
\begin{itemize}
    \item Finish = [false, false, false, false, false]
    \item Work = [2, 3, 0]
\end{itemize}

\noindent \textbf{Thực hiện thuật toán an toàn:}
\begin{itemize}
    \item $P_1$ có Need <= Work, nên Work = Work + Allocation của $P_1$ = [5, 3, 2], Finish[1] = true
    \item $P_3$ có Need <= Work, nên Work = Work + Allocation của $P_3$ = [7, 4, 3], Finish[3] = true
    \item $P_4$ có Need <= Work, nên Work = Work + Allocation của $P_4$ = [7, 4, 5], Finish[4] = true
    \item $P_0$ có Need <= Work, nên Work = Work + Allocation của $P_0$ = [7, 5, 5], Finish[0] = true
    \item $P_2$ có Need <= Work, nên Work = Work + Allocation của $P_2$ = [10, 5, 7], Finish[2] = true
\end{itemize}

\noindent \textbf{Kết luận:} Chuỗi < $P_1, P_3, P_4, P_0, P_2$ > thoả mãn điều kiện an toàn vì tất cả Finish[i] = true.

\subsection{Kiểm tra yêu cầu $P_4$}
Vì Need của $P_4$ (4, 3, 1) <= Available (2, 3, 0), yêu cầu (3, 3, 0) của $P_4$ không được chấp nhận.

\subsection{Kiểm tra yêu cầu $P_0$}
Vì Need của $P_0$ (7, 4, 3) <= Available (2, 3, 0), yêu cầu (0, 2, 0) của $P_0$ không được chấp nhận.
\section{Bài tập slide trang 34, 35}
\subsection{Nội dung bài tập}
\begin{table}[H]
    \centering
    \begin{tabular}{|c|c|c|c|}
        \toprule
        & \textbf{Allocation} & \textbf{Request} & \textbf{Available} \\
        & A B C & A B C & A B C \\
        \midrule
        $P_0$ & 0 1 0 & 0 0 0 & 0 0 0 \\
        $P_1$ & 2 0 0 & 2 0 2 &\\
        $P_2$ & 3 0 3 & 0 0 0 &\\
        $P_3$ & 2 1 1 & 1 0 0 &\\
        $P_4$ & 0 0 2 & 0 0 2 &\\
        \bottomrule
    \end{tabular}
\end{table}

\begin{itemize}
    \item Chuỗi < $P_0, P_2, P_3, P_1, P_4$ > cho phép $Finish[i]$ = true với tất cả $0 <= i <= n - 1$ 
    \item $P_2$ yêu cầu thêm một đơn vị tài nguyên C
\end{itemize}

\begin{table}[H]
    \centering
    \begin{tabular}{|c|c|}
        \toprule
        & \textbf{Request}\\
        & A B C\\
        \midrule
        $P_0$ & 0 0 0 \\
        $P_1$ & 2 0 2 \\
        $P_2$ & 0 0 1 \\
        $P_3$ & 1 0 0 \\
        $P_4$ & 0 0 2 \\
        \bottomrule
    \end{tabular}
\end{table}
\begin{itemize}
    \item Hệ thống đang ở trạng thái nào?
\end{itemize}

\subsection{Giải bài tập}
\subsection{Khởi tạo 2 mảng Finish và Work}
\begin{itemize}
    \item Finish = [false, false, false, false, false]
    \item Work = [0, 0, 0]
\end{itemize}

\noindent \textbf{Thực hiện thuật toán an toàn:}
\begin{itemize}
    \item $P_0$ có Request <= Work, nên Work = Work + Allocation của $P_0$ = [0, 1, 0], Finish[0] = true
    \item $P_2$ có Request <= Work, nên Work = Work + Allocation của $P_2$ = [3, 1, 3], Finish[2] = true
    \item $P_3$ có Request <= Work, nên Work = Work + Allocation của $P_3$ = [5, 2, 4], Finish[3] = true
    \item $P_1$ có Request <= Work, nên Work = Work + Allocation của $P_1$ = [7, 2, 4], Finish[1] = true
    \item $P_4$ có Request <= Work, nên Work = Work + Allocation của $P_4$ = [7, 2, 6], Finish[4] = true
\end{itemize}

\noindent \textbf{Kết luận:} Chuỗi < $P_0, P_2, P_3, P_1, P_4$ > cho phép $Finish[i]$ = true với tất cả $0 <= i <= n - 1$, do đó hệ thống đang ở trạng thái an toàn.

\subsection{Kiểm tra yêu cầu $P_2$}
\noindent \textbf{Thực hiện thuật toán an toàn:}
\begin{itemize}
    \item $P_0$ có Request <= Work, nên Work = Work + Allocation của $P_0$ = [0, 1, 0], Finish[0] = true
    \item $P_2$ có Request > Work, nên yêu cầu của $P_2$ không được chấp nhận, Finish[2] = false
    \item $P_3$ có Request > Work, nên yêu cầu của $P_3$ không được chấp nhận, Finish[3] = false
    \item $P_1$ có Request > Work, nên yêu cầu của $P_1$ không được chấp nhận, Finish[1] = false
    \item $P_4$ có Request > Work, nên yêu cầu của $P_4$ không được chấp nhận, Finish[4] = false
\end{itemize}

\noindent \textbf{Kết luận:} Yêu cầu của $P_2$ không được chấp nhận vì nó vượt quá tài nguyên hiện có, do đó hệ thống đang ở trạng thái không an toàn.
\end{document}